% Generated draft: requires mathematical and source review.
\documentclass[11pt]{article}
\usepackage[margin=25mm]{geometry}
\usepackage{fontspec}
\setmainfont{Latin Modern Roman}
\usepackage{amsmath,amssymb,mathtools,unicode-math}
\setmathfont{Latin Modern Math}
\usepackage{enumitem,needspace,xurl,hyperref}
\setlength{\parindent}{0pt}
\setlength{\parskip}{5pt}
\setlength{\emergencystretch}{3em}
\newcommand{\N}{\mathbb{N}}
\newcommand{\Z}{\mathbb{Z}}
\newcommand{\Q}{\mathbb{Q}}
\newcommand{\R}{\mathbb{R}}
\newcommand{\C}{\mathbb{C}}
\newcommand{\F}{\mathbb{F}}
\newcommand{\A}{\mathbb{A}}
\newcommand{\PP}{\mathbb{P}}
\newcommand{\E}{\mathbb{E}}
\newcommand{\eps}{\varepsilon}
\newcommand{\dd}{\,\mathrm{d}}
\begin{document}
\section*{UnsolvedMath ID: 11000165}
\textbf{Source problem code:} catalogue entry
\subsection{Definitions and mathematical statement}
In the usual closed-surface conventions, let \(\operatorname{Mod}_g\) be the genus-\(g\) mapping class group, \(H_g\) its subgroup of classes extending over a fixed marked genus-\(g\) handlebody, and \(I_g\) its Torelli subgroup, the kernel of the action on \(H_1(\Sigma_g;\Z)\). The notation \(K_g\) conventionally denotes the Johnson kernel, but its precise definition and low-genus conventions are not supplied here.

Problem 2.8 on PDF page 157 proposes finding generators for
\[
 H_g\cap I_g\qquad\text{and/or}\qquad H_g\cap K_g.
\]
A generating family must generate the intersection as a group, with the source's boundary and marking conditions. The record does not require the family to be finite and does not replace generation by normal generation. It gives no proposed elements, genus range or uniformity criterion.

The source motivates this by its earlier Heegaard parametrization: it reports that every genus-\(g\) Heegaard splitting of a homology \(3\)-sphere is obtained by ranging a parameter \(\phi\) over \(I_g\), and that Morita [36] permits ranging it over \(K_g\). The fixed reference gluing and parametrization are absent. In particular this must not be interpreted as gluing two identically marked handlebodies by \(\phi\) alone without the required reference data. Homology-sphere coefficients and the equivalence of splittings also need confirmation. The last sentence introduces a different subsequent problem about graph-manifold splittings; its actual statement is not included.
\subsection{Short English statement}
Find generating families for the handlebody group's intersections with the Torelli group and/or Johnson kernel. The source motivates this through Heegaard parametrizations of homology spheres, but omits the reference gluing and subgroup conventions. No finite-generation requirement is stated.
\subsection{Sources}
\begin{itemize}
\item Local UnsolvedMath registry, permanent ID \(11000165\), exact source code \texttt{catalogue entry}, Problem 2.8, PDF page 157. No individual proposer or proposed year is supplied; the title is truncated.
\item Farb, edited volume \emph{Problems on Mapping Class Groups and Related Topics}: \url{https://www.math.uchicago.edu/~farb/papers/mcgbook.pdf}. The volume has 329 explicitly typeset numbered items across chapters, repeated local numbering distinguished by PDF page.
\item Morita [36] is cited for the \(K_g\) Heegaard-parametrization result; its complete bibliography entry and the preceding reference construction are absent.
\item Extraction pdf-font-structure, accessed 2026-07-29. Source evidence NEEDS\_REVIEW presents an open problem and requires a current primary-literature audit. Registry status open and classification OPEN-TRIAGE; the research summary says only ``Literature status: Truncated.''
\item Dataset attribution: Ulam AI / UnsolvedMath Contributors, CC BY 4.0. Source rights NEEDS\_REVIEW: publicly accessible author/publisher copy, redistribution terms requiring review.
\end{itemize}
\end{document}
